% =====================================================================
\section{Algorithm settings}
\label{app:algorithm}
% =====================================================================
% Needs \appendix before % =====================================================================
\section{Algorithm settings}
\label{app:algorithm}
% =====================================================================
% Needs \appendix before % =====================================================================
\section{Algorithm settings}
\label{app:algorithm}
% =====================================================================
% Needs \appendix before % =====================================================================
\section{Algorithm settings}
\label{app:algorithm}
% =====================================================================
% Needs \appendix before \input{sections/appendix_algorithm} in nvision_paper.tex, after the Conclusion.
% Every number below is the code default (nvision/belief/smc_marginal.py, nvision/sim/defaults.py,
% nvision/sim/locs/bayesian/sbed_locator.py, nvision/belief/focus_window.py, nvision/belief/dip_detection.py,
% nvision/sim/locs/search_grid.py).

\paragraph{Particle filter.} The cloud has $N_p=1000$ particles. Every dimension is scaled to $[0,1]$ by its prior
bounds, so one resampling kernel serves all of them; reported values are mapped back to Hz. The cloud is resampled
when the effective sample size $(\sum_i w_i^2)^{-1}$ falls below $0.2\,N_p$, systematically, and then rejuvenated with
the Liu--West kernel,
\begin{equation}
  \psi_i'\sim\mathcal N\big(a\psi_i+(1-a)\bar\psi,\ (1-a^2)\Sigma\big),\qquad a=0.98,
  \label{eq:liuwest}
\end{equation}
with $\psi_i=(\theta_i,\phi_i)$ the whole particle and $\bar\psi$, $\Sigma$ the weighted mean and covariance before
resampling. The diagonal of $\Sigma$ is floored at $(0.01\,e^{-n/25})^2$ of each unit range after $n$ measurements,
and a particle that leaves the unit cube is reflected back into it.

\paragraph{Noise model.} Each particle carries an inverse-gamma posterior on $\sigma^2$ with shape $\alpha$, the same
for every particle because its update does not depend on the residuals, and scale $\beta_i$. The prior is
$\alpha_0=10$, $\beta_0=\alpha_0\sigma_{\rm nom}^2$, with $\sigma_{\rm nom}$ the geometric mean of the prior range of
$\sigma$. Before each measurement both are discounted by $d=0.99$; a measurement $y_n$, the mean of $k$ shots at
$x_n$, then adds $\tfrac12$ to $\alpha$ and $\tfrac12 k\,r_{n,i}^2$ to $\beta_i$, with residual
$r_{n,i}=y_n-S(x_n;\theta_i,\phi_i)$, so that after $n$ measurements
$\beta_i=d^n\beta_0+\tfrac12\sum_{t<n}d^{\,n-1-t}k\,r_{t,i}^2$. Old residuals are thus forgotten with a memory of
about $1/(1-d)=100$ measurements, and $\alpha$ saturates at $0.5/(1-d)=50$, about 100 degrees of freedom. With
$\sigma^2$ integrated out, the measurement reweights the particles by the Student-$t$ predictive
\begin{equation}
  w_i\ \propto\ w_i\;t_{2\alpha}\!\Big(y_n;\,S(x_n;\theta_i,\phi_i),\,\frac{\beta_i}{k\alpha}\Big),
  \label{eq:smc-update}
\end{equation}
taken with $\alpha$, $\beta_i$ before the update. After a rejuvenation each $\beta_i$ is recomputed from the full
history at the particle's new position.

\paragraph{Initial cloud.} \tbd{The prior on $\Delta_Z$ is being changed to uniform over the range the truth is drawn
from (now a truncated Gaussian, offset from the truth, plus 10\% of the particles uniform over the bounds); state it
here once Sec.~\ref{sec:prior} is fixed.}

\paragraph{Stop rule and budget.} The ceilings on the claimed uncertainty are $100$~kHz for $\theta$, $500$~kHz for
each width and $0.01$ for the contrast; they must hold for 8 consecutive measurements, together with a dip found by
the dip test, which is re-run at most every 10 measurements. Every strategy has the same budget of 2500
measurements; a sweep's total is capped in whole passes. A run that reaches the budget without stopping is censored
and counted at the budget, as a lower bound on its cost.

\paragraph{Dip test.} Each measurement gets the score $z_n=(b-y_n)/s$. The baseline $b$ is the 70th percentile of
the $y_n$ minus its expected offset under pure noise, $t_\nu^{-1}(0.7)\,s$, and $s$ and $\nu$ are the scale and
degrees of freedom of the noise posterior. Every contiguous stretch of measured drive frequencies whose span lies in
the prior range of the linewidth gets the boxcar score $\sum z_n/\sqrt{n}$, over its $n$ measurements. A stretch is
significant when its score exceeds $t_\nu^{-1}(1-p_{\rm FA}/N_{\rm tests})$, the Bonferroni threshold for a
family-wise false-alarm probability $p_{\rm FA}=0.1$ over the $N_{\rm tests}$ stretches tested.

\paragraph{Sweep family.} $\Gamma_0$ is the FWHM of the prior-mean line. SimpleSweep measures a uniform grid of
spacing $\Gamma_0/p$, $p=2,3,4$. Smart-Sweep sweeps at $\Gamma_0/3$ until the dip test finds a dip, then at
$\Gamma_0/10$ inside the focus, and then measures only at the maxima of $g(x)^2$ at the posterior mean, cycling
between them, once the posterior expects each of those points to carry at least $90\%$ of the information of the
true optimum.

\paragraph{Bayesian SBED.} After every resampling SBED places 800 candidate drive frequencies at the quantiles of a
density with Gaussian kernels on the flanks, $\pm$HWHM, of the posterior-mean line, kernels on the dips found and a
flat share, thinned to a spacing of at least $100$~kHz. Eq.~\eqref{eq:eig-var} is evaluated on a 500-particle
subsample, with $\sigma^2$ the posterior mean of the noise. The candidates are taken in blocks of 64, and the block
maxima compete through a softmax with a temperature of $0.01$~nats, so the choice is near-greedy. Each focus section
is a dip's extent padded by one posterior-mean HWHM on each side. A new dip widens the focus at once; narrowing waits
8 measurements and must shrink the focus by at least 5\%. The focus is re-decided after every resampling. The
falsification probe is made with probability $0.02+0.08\,g$, where $g$ is the widest unmeasured gap as a fraction of
the window.

\paragraph{Use cases.} The simulated signals are pseudo-Voigt dips of total FWHM 2~MHz, on a grid of the SNR, the
window (1.6, 2, 3 and 5 times the width one dip group covers), the Gaussian share of the width (0, 0.15, 0.3) and the
isotope (unresolved hyperfine lines, $^{14}$N or $^{15}$N). The code labels the SNR levels 1, 2, 5, 10, 20 against a
contrast of 0.25 \tbd{convert to $C/\sigma$ of Eq.~\eqref{eq:snr}: the code's contrast is $2C$}. \tbd{Number of
repeats per use case.}
 in nvision_paper.tex, after the Conclusion.
% Every number below is the code default (nvision/belief/smc_marginal.py, nvision/sim/defaults.py,
% nvision/sim/locs/bayesian/sbed_locator.py, nvision/belief/focus_window.py, nvision/belief/dip_detection.py,
% nvision/sim/locs/search_grid.py).

\paragraph{Particle filter.} The cloud has $N_p=1000$ particles. Every dimension is scaled to $[0,1]$ by its prior
bounds, so one resampling kernel serves all of them; reported values are mapped back to Hz. The cloud is resampled
when the effective sample size $(\sum_i w_i^2)^{-1}$ falls below $0.2\,N_p$, systematically, and then rejuvenated with
the Liu--West kernel,
\begin{equation}
  \psi_i'\sim\mathcal N\big(a\psi_i+(1-a)\bar\psi,\ (1-a^2)\Sigma\big),\qquad a=0.98,
  \label{eq:liuwest}
\end{equation}
with $\psi_i=(\theta_i,\phi_i)$ the whole particle and $\bar\psi$, $\Sigma$ the weighted mean and covariance before
resampling. The diagonal of $\Sigma$ is floored at $(0.01\,e^{-n/25})^2$ of each unit range after $n$ measurements,
and a particle that leaves the unit cube is reflected back into it.

\paragraph{Noise model.} Each particle carries an inverse-gamma posterior on $\sigma^2$ with shape $\alpha$, the same
for every particle because its update does not depend on the residuals, and scale $\beta_i$. The prior is
$\alpha_0=10$, $\beta_0=\alpha_0\sigma_{\rm nom}^2$, with $\sigma_{\rm nom}$ the geometric mean of the prior range of
$\sigma$. Before each measurement both are discounted by $d=0.99$; a measurement $y_n$, the mean of $k$ shots at
$x_n$, then adds $\tfrac12$ to $\alpha$ and $\tfrac12 k\,r_{n,i}^2$ to $\beta_i$, with residual
$r_{n,i}=y_n-S(x_n;\theta_i,\phi_i)$, so that after $n$ measurements
$\beta_i=d^n\beta_0+\tfrac12\sum_{t<n}d^{\,n-1-t}k\,r_{t,i}^2$. Old residuals are thus forgotten with a memory of
about $1/(1-d)=100$ measurements, and $\alpha$ saturates at $0.5/(1-d)=50$, about 100 degrees of freedom. With
$\sigma^2$ integrated out, the measurement reweights the particles by the Student-$t$ predictive
\begin{equation}
  w_i\ \propto\ w_i\;t_{2\alpha}\!\Big(y_n;\,S(x_n;\theta_i,\phi_i),\,\frac{\beta_i}{k\alpha}\Big),
  \label{eq:smc-update}
\end{equation}
taken with $\alpha$, $\beta_i$ before the update. After a rejuvenation each $\beta_i$ is recomputed from the full
history at the particle's new position.

\paragraph{Initial cloud.} \tbd{The prior on $\Delta_Z$ is being changed to uniform over the range the truth is drawn
from (now a truncated Gaussian, offset from the truth, plus 10\% of the particles uniform over the bounds); state it
here once Sec.~\ref{sec:prior} is fixed.}

\paragraph{Stop rule and budget.} The ceilings on the claimed uncertainty are $100$~kHz for $\theta$, $500$~kHz for
each width and $0.01$ for the contrast; they must hold for 8 consecutive measurements, together with a dip found by
the dip test, which is re-run at most every 10 measurements. Every strategy has the same budget of 2500
measurements; a sweep's total is capped in whole passes. A run that reaches the budget without stopping is censored
and counted at the budget, as a lower bound on its cost.

\paragraph{Dip test.} Each measurement gets the score $z_n=(b-y_n)/s$. The baseline $b$ is the 70th percentile of
the $y_n$ minus its expected offset under pure noise, $t_\nu^{-1}(0.7)\,s$, and $s$ and $\nu$ are the scale and
degrees of freedom of the noise posterior. Every contiguous stretch of measured drive frequencies whose span lies in
the prior range of the linewidth gets the boxcar score $\sum z_n/\sqrt{n}$, over its $n$ measurements. A stretch is
significant when its score exceeds $t_\nu^{-1}(1-p_{\rm FA}/N_{\rm tests})$, the Bonferroni threshold for a
family-wise false-alarm probability $p_{\rm FA}=0.1$ over the $N_{\rm tests}$ stretches tested.

\paragraph{Sweep family.} $\Gamma_0$ is the FWHM of the prior-mean line. SimpleSweep measures a uniform grid of
spacing $\Gamma_0/p$, $p=2,3,4$. Smart-Sweep sweeps at $\Gamma_0/3$ until the dip test finds a dip, then at
$\Gamma_0/10$ inside the focus, and then measures only at the maxima of $g(x)^2$ at the posterior mean, cycling
between them, once the posterior expects each of those points to carry at least $90\%$ of the information of the
true optimum.

\paragraph{Bayesian SBED.} After every resampling SBED places 800 candidate drive frequencies at the quantiles of a
density with Gaussian kernels on the flanks, $\pm$HWHM, of the posterior-mean line, kernels on the dips found and a
flat share, thinned to a spacing of at least $100$~kHz. Eq.~\eqref{eq:eig-var} is evaluated on a 500-particle
subsample, with $\sigma^2$ the posterior mean of the noise. The candidates are taken in blocks of 64, and the block
maxima compete through a softmax with a temperature of $0.01$~nats, so the choice is near-greedy. Each focus section
is a dip's extent padded by one posterior-mean HWHM on each side. A new dip widens the focus at once; narrowing waits
8 measurements and must shrink the focus by at least 5\%. The focus is re-decided after every resampling. The
falsification probe is made with probability $0.02+0.08\,g$, where $g$ is the widest unmeasured gap as a fraction of
the window.

\paragraph{Use cases.} The simulated signals are pseudo-Voigt dips of total FWHM 2~MHz, on a grid of the SNR, the
window (1.6, 2, 3 and 5 times the width one dip group covers), the Gaussian share of the width (0, 0.15, 0.3) and the
isotope (unresolved hyperfine lines, $^{14}$N or $^{15}$N). The code labels the SNR levels 1, 2, 5, 10, 20 against a
contrast of 0.25 \tbd{convert to $C/\sigma$ of Eq.~\eqref{eq:snr}: the code's contrast is $2C$}. \tbd{Number of
repeats per use case.}
 in nvision_paper.tex, after the Conclusion.
% Every number below is the code default (nvision/belief/smc_marginal.py, nvision/sim/defaults.py,
% nvision/sim/locs/bayesian/sbed_locator.py, nvision/belief/focus_window.py, nvision/belief/dip_detection.py,
% nvision/sim/locs/search_grid.py).

\paragraph{Particle filter.} The cloud has $N_p=1000$ particles. Every dimension is scaled to $[0,1]$ by its prior
bounds, so one resampling kernel serves all of them; reported values are mapped back to Hz. The cloud is resampled
when the effective sample size $(\sum_i w_i^2)^{-1}$ falls below $0.2\,N_p$, systematically, and then rejuvenated with
the Liu--West kernel,
\begin{equation}
  \psi_i'\sim\mathcal N\big(a\psi_i+(1-a)\bar\psi,\ (1-a^2)\Sigma\big),\qquad a=0.98,
  \label{eq:liuwest}
\end{equation}
with $\psi_i=(\theta_i,\phi_i)$ the whole particle and $\bar\psi$, $\Sigma$ the weighted mean and covariance before
resampling. The diagonal of $\Sigma$ is floored at $(0.01\,e^{-n/25})^2$ of each unit range after $n$ measurements,
and a particle that leaves the unit cube is reflected back into it.

\paragraph{Noise model.} Each particle carries an inverse-gamma posterior on $\sigma^2$ with shape $\alpha$, the same
for every particle because its update does not depend on the residuals, and scale $\beta_i$. The prior is
$\alpha_0=10$, $\beta_0=\alpha_0\sigma_{\rm nom}^2$, with $\sigma_{\rm nom}$ the geometric mean of the prior range of
$\sigma$. Before each measurement both are discounted by $d=0.99$; a measurement $y_n$, the mean of $k$ shots at
$x_n$, then adds $\tfrac12$ to $\alpha$ and $\tfrac12 k\,r_{n,i}^2$ to $\beta_i$, with residual
$r_{n,i}=y_n-S(x_n;\theta_i,\phi_i)$, so that after $n$ measurements
$\beta_i=d^n\beta_0+\tfrac12\sum_{t<n}d^{\,n-1-t}k\,r_{t,i}^2$. Old residuals are thus forgotten with a memory of
about $1/(1-d)=100$ measurements, and $\alpha$ saturates at $0.5/(1-d)=50$, about 100 degrees of freedom. With
$\sigma^2$ integrated out, the measurement reweights the particles by the Student-$t$ predictive
\begin{equation}
  w_i\ \propto\ w_i\;t_{2\alpha}\!\Big(y_n;\,S(x_n;\theta_i,\phi_i),\,\frac{\beta_i}{k\alpha}\Big),
  \label{eq:smc-update}
\end{equation}
taken with $\alpha$, $\beta_i$ before the update. After a rejuvenation each $\beta_i$ is recomputed from the full
history at the particle's new position.

\paragraph{Initial cloud.} \tbd{The prior on $\Delta_Z$ is being changed to uniform over the range the truth is drawn
from (now a truncated Gaussian, offset from the truth, plus 10\% of the particles uniform over the bounds); state it
here once Sec.~\ref{sec:prior} is fixed.}

\paragraph{Stop rule and budget.} The ceilings on the claimed uncertainty are $100$~kHz for $\theta$, $500$~kHz for
each width and $0.01$ for the contrast; they must hold for 8 consecutive measurements, together with a dip found by
the dip test, which is re-run at most every 10 measurements. Every strategy has the same budget of 2500
measurements; a sweep's total is capped in whole passes. A run that reaches the budget without stopping is censored
and counted at the budget, as a lower bound on its cost.

\paragraph{Dip test.} Each measurement gets the score $z_n=(b-y_n)/s$. The baseline $b$ is the 70th percentile of
the $y_n$ minus its expected offset under pure noise, $t_\nu^{-1}(0.7)\,s$, and $s$ and $\nu$ are the scale and
degrees of freedom of the noise posterior. Every contiguous stretch of measured drive frequencies whose span lies in
the prior range of the linewidth gets the boxcar score $\sum z_n/\sqrt{n}$, over its $n$ measurements. A stretch is
significant when its score exceeds $t_\nu^{-1}(1-p_{\rm FA}/N_{\rm tests})$, the Bonferroni threshold for a
family-wise false-alarm probability $p_{\rm FA}=0.1$ over the $N_{\rm tests}$ stretches tested.

\paragraph{Sweep family.} $\Gamma_0$ is the FWHM of the prior-mean line. SimpleSweep measures a uniform grid of
spacing $\Gamma_0/p$, $p=2,3,4$. Smart-Sweep sweeps at $\Gamma_0/3$ until the dip test finds a dip, then at
$\Gamma_0/10$ inside the focus, and then measures only at the maxima of $g(x)^2$ at the posterior mean, cycling
between them, once the posterior expects each of those points to carry at least $90\%$ of the information of the
true optimum.

\paragraph{Bayesian SBED.} After every resampling SBED places 800 candidate drive frequencies at the quantiles of a
density with Gaussian kernels on the flanks, $\pm$HWHM, of the posterior-mean line, kernels on the dips found and a
flat share, thinned to a spacing of at least $100$~kHz. Eq.~\eqref{eq:eig-var} is evaluated on a 500-particle
subsample, with $\sigma^2$ the posterior mean of the noise. The candidates are taken in blocks of 64, and the block
maxima compete through a softmax with a temperature of $0.01$~nats, so the choice is near-greedy. Each focus section
is a dip's extent padded by one posterior-mean HWHM on each side. A new dip widens the focus at once; narrowing waits
8 measurements and must shrink the focus by at least 5\%. The focus is re-decided after every resampling. The
falsification probe is made with probability $0.02+0.08\,g$, where $g$ is the widest unmeasured gap as a fraction of
the window.

\paragraph{Use cases.} The simulated signals are pseudo-Voigt dips of total FWHM 2~MHz, on a grid of the SNR, the
window (1.6, 2, 3 and 5 times the width one dip group covers), the Gaussian share of the width (0, 0.15, 0.3) and the
isotope (unresolved hyperfine lines, $^{14}$N or $^{15}$N). The code labels the SNR levels 1, 2, 5, 10, 20 against a
contrast of 0.25 \tbd{convert to $C/\sigma$ of Eq.~\eqref{eq:snr}: the code's contrast is $2C$}. \tbd{Number of
repeats per use case.}
 in nvision_paper.tex, after the Conclusion.
% Every number below is the code default (nvision/belief/smc_marginal.py, nvision/sim/defaults.py,
% nvision/sim/locs/bayesian/sbed_locator.py, nvision/belief/focus_window.py, nvision/belief/dip_detection.py,
% nvision/sim/locs/search_grid.py).

\paragraph{Particle filter.} The cloud has $N_p=1000$ particles. Every dimension is scaled to $[0,1]$ by its prior
bounds, so one resampling kernel serves all of them; reported values are mapped back to Hz. The cloud is resampled
when the effective sample size $(\sum_i w_i^2)^{-1}$ falls below $0.2\,N_p$, systematically, and then rejuvenated with
the Liu--West kernel,
\begin{equation}
  \psi_i'\sim\mathcal N\big(a\psi_i+(1-a)\bar\psi,\ (1-a^2)\Sigma\big),\qquad a=0.98,
  \label{eq:liuwest}
\end{equation}
with $\psi_i=(\theta_i,\phi_i)$ the whole particle and $\bar\psi$, $\Sigma$ the weighted mean and covariance before
resampling. The diagonal of $\Sigma$ is floored at $(0.01\,e^{-n/25})^2$ of each unit range after $n$ measurements,
and a particle that leaves the unit cube is reflected back into it.

\paragraph{Noise model.} Each particle carries an inverse-gamma posterior on $\sigma^2$ with shape $\alpha$, the same
for every particle because its update does not depend on the residuals, and scale $\beta_i$. The prior is
$\alpha_0=10$, $\beta_0=\alpha_0\sigma_{\rm nom}^2$, with $\sigma_{\rm nom}$ the geometric mean of the prior range of
$\sigma$. Before each measurement both are discounted by $d=0.99$; a measurement $y_n$, the mean of $k$ shots at
$x_n$, then adds $\tfrac12$ to $\alpha$ and $\tfrac12 k\,r_{n,i}^2$ to $\beta_i$, with residual
$r_{n,i}=y_n-S(x_n;\theta_i,\phi_i)$, so that after $n$ measurements
$\beta_i=d^n\beta_0+\tfrac12\sum_{t<n}d^{\,n-1-t}k\,r_{t,i}^2$. Old residuals are thus forgotten with a memory of
about $1/(1-d)=100$ measurements, and $\alpha$ saturates at $0.5/(1-d)=50$, about 100 degrees of freedom. With
$\sigma^2$ integrated out, the measurement reweights the particles by the Student-$t$ predictive
\begin{equation}
  w_i\ \propto\ w_i\;t_{2\alpha}\!\Big(y_n;\,S(x_n;\theta_i,\phi_i),\,\frac{\beta_i}{k\alpha}\Big),
  \label{eq:smc-update}
\end{equation}
taken with $\alpha$, $\beta_i$ before the update. After a rejuvenation each $\beta_i$ is recomputed from the full
history at the particle's new position.

\paragraph{Initial cloud.} \tbd{The prior on $\Delta_Z$ is being changed to uniform over the range the truth is drawn
from (now a truncated Gaussian, offset from the truth, plus 10\% of the particles uniform over the bounds); state it
here once Sec.~\ref{sec:prior} is fixed.}

\paragraph{Stop rule and budget.} The ceilings on the claimed uncertainty are $100$~kHz for $\theta$, $500$~kHz for
each width and $0.01$ for the contrast; they must hold for 8 consecutive measurements, together with a dip found by
the dip test, which is re-run at most every 10 measurements. Every strategy has the same budget of 2500
measurements; a sweep's total is capped in whole passes. A run that reaches the budget without stopping is censored
and counted at the budget, as a lower bound on its cost.

\paragraph{Dip test.} Each measurement gets the score $z_n=(b-y_n)/s$. The baseline $b$ is the 70th percentile of
the $y_n$ minus its expected offset under pure noise, $t_\nu^{-1}(0.7)\,s$, and $s$ and $\nu$ are the scale and
degrees of freedom of the noise posterior. Every contiguous stretch of measured drive frequencies whose span lies in
the prior range of the linewidth gets the boxcar score $\sum z_n/\sqrt{n}$, over its $n$ measurements. A stretch is
significant when its score exceeds $t_\nu^{-1}(1-p_{\rm FA}/N_{\rm tests})$, the Bonferroni threshold for a
family-wise false-alarm probability $p_{\rm FA}=0.1$ over the $N_{\rm tests}$ stretches tested.

\paragraph{Sweep family.} $\Gamma_0$ is the FWHM of the prior-mean line. SimpleSweep measures a uniform grid of
spacing $\Gamma_0/p$, $p=2,3,4$. Smart-Sweep sweeps at $\Gamma_0/3$ until the dip test finds a dip, then at
$\Gamma_0/10$ inside the focus, and then measures only at the maxima of $g(x)^2$ at the posterior mean, cycling
between them, once the posterior expects each of those points to carry at least $90\%$ of the information of the
true optimum.

\paragraph{Bayesian SBED.} After every resampling SBED places 800 candidate drive frequencies at the quantiles of a
density with Gaussian kernels on the flanks, $\pm$HWHM, of the posterior-mean line, kernels on the dips found and a
flat share, thinned to a spacing of at least $100$~kHz. Eq.~\eqref{eq:eig-var} is evaluated on a 500-particle
subsample, with $\sigma^2$ the posterior mean of the noise. The candidates are taken in blocks of 64, and the block
maxima compete through a softmax with a temperature of $0.01$~nats, so the choice is near-greedy. Each focus section
is a dip's extent padded by one posterior-mean HWHM on each side. A new dip widens the focus at once; narrowing waits
8 measurements and must shrink the focus by at least 5\%. The focus is re-decided after every resampling. The
falsification probe is made with probability $0.02+0.08\,g$, where $g$ is the widest unmeasured gap as a fraction of
the window.

\paragraph{Use cases.} The simulated signals are pseudo-Voigt dips of total FWHM 2~MHz, on a grid of the SNR, the
window (1.6, 2, 3 and 5 times the width one dip group covers), the Gaussian share of the width (0, 0.15, 0.3) and the
isotope (unresolved hyperfine lines, $^{14}$N or $^{15}$N). The code labels the SNR levels 1, 2, 5, 10, 20 against a
contrast of 0.25 \tbd{convert to $C/\sigma$ of Eq.~\eqref{eq:snr}: the code's contrast is $2C$}. \tbd{Number of
repeats per use case.}
